% Options for packages loaded elsewhere
\PassOptionsToPackage{unicode}{hyperref}
\PassOptionsToPackage{hyphens}{url}
\documentclass[
  11pt,
]{article}
\usepackage{xcolor}
\usepackage[margin=1in]{geometry}
\usepackage{amsmath,amssymb}
\setcounter{secnumdepth}{-\maxdimen} % remove section numbering
\usepackage{iftex}
\ifPDFTeX
  \usepackage[T1]{fontenc}
  \usepackage[utf8]{inputenc}
  \usepackage{textcomp} % provide euro and other symbols
\else % if luatex or xetex
  \usepackage{unicode-math} % this also loads fontspec
  \defaultfontfeatures{Scale=MatchLowercase}
  \defaultfontfeatures[\rmfamily]{Ligatures=TeX,Scale=1}
\fi
\usepackage{lmodern}
\ifPDFTeX\else
  % xetex/luatex font selection
\fi
% Use upquote if available, for straight quotes in verbatim environments
\IfFileExists{upquote.sty}{\usepackage{upquote}}{}
\IfFileExists{microtype.sty}{% use microtype if available
  \usepackage[]{microtype}
  \UseMicrotypeSet[protrusion]{basicmath} % disable protrusion for tt fonts
}{}
\usepackage{setspace}
\makeatletter
\@ifundefined{KOMAClassName}{% if non-KOMA class
  \IfFileExists{parskip.sty}{%
    \usepackage{parskip}
  }{% else
    \setlength{\parindent}{0pt}
    \setlength{\parskip}{6pt plus 2pt minus 1pt}}
}{% if KOMA class
  \KOMAoptions{parskip=half}}
\makeatother
\setlength{\emergencystretch}{3em} % prevent overfull lines
\providecommand{\tightlist}{%
  \setlength{\itemsep}{0pt}\setlength{\parskip}{0pt}}
\usepackage{bookmark}
\IfFileExists{xurl.sty}{\usepackage{xurl}}{} % add URL line breaks if available
\urlstyle{same}
\hypersetup{
  pdftitle={Quantum Mechanics as a Forced Structural Consequence of Event Density Primitives},
  pdfauthor={Allen Proxmire},
  hidelinks,
  pdfcreator={LaTeX via pandoc}}

\title{Quantum Mechanics as a Forced Structural Consequence of Event
Density Primitives}
\author{Allen Proxmire}
\date{April 2026}

\begin{document}
\maketitle

\setstretch{1.15}
\section{Quantum Mechanics as a Forced Structural Consequence of Event
Density
Primitives}\label{quantum-mechanics-as-a-forced-structural-consequence-of-event-density-primitives}

\textbf{Allen Proxmire}

\textbf{April 2026}

\begin{center}\rule{0.5\linewidth}{0.5pt}\end{center}

\subsection{Abstract}\label{abstract}

The non-relativistic quantum mechanics of a single particle is shown to
arise as a structural consequence of the Event Density (ED) primitive
stack. A single mathematical object --- the participation measure
\(P_K(x,t) = \sqrt{b_K(x,t)} \, e^{i \pi_K(x,t)}\), a complex-valued
distribution over channels \(K\) and positions \(x\) --- fuses two ED
primitives (bandwidth and tension polarity) into the carrier of
quantum-mechanical content. In a well-defined thin-participation limit,
the coherent sum \(\Psi = \sum_K P_K\) satisfies the standard linear
Schrödinger equation. Environmental phase-randomization at commitment
events, together with an individuation threshold, forces the Born rule
in its squared-amplitude form as a definitional consequence of
\(b_K = |P_K|^2\). Non-factorizable joint participation measures on
bipartite systems produce Bell-inequality-violating correlations that
saturate the Tsirelson bound \(2\sqrt{2}\). The Fourier-conjugate
decomposition of the adjacency-band bandwidth budget, combined with the
Fourier-uncertainty theorem, yields the Heisenberg uncertainty relation
\(\Delta x \cdot \Delta p \geq \hbar/2\). A single structural commitment
--- the identification of bandwidth with squared amplitude,
\(b_K = |P_K|^2\) --- threads through five otherwise-independent
appearances of the number two in quantum mechanics: the Born rule, the
Schrödinger kinetic term, the Madelung amplitude-phase decomposition,
the sublinear bandwidth-composition rule, and the Fourier-uncertainty
bound. The program derives the form of quantum-mechanical equations; the
numerical values of the dimensional constants \(\hbar\) and \(m\) are
inherited from the Event Density Dimensional Atlas rather than derived
from primitives, and the primitive-level derivation of particle masses
is flagged as a frontier research problem. No physical postulate beyond
the participation-measure structure, the four-band bandwidth
decomposition, and the commitment dynamics is required. The result is a
structural reduction of the four independent postulates of standard
quantum mechanics to consequences of a single primitive-level framework.

\begin{center}\rule{0.5\linewidth}{0.5pt}\end{center}

\subsection{1. Introduction}\label{introduction}

Standard non-relativistic quantum mechanics rests on four logically
independent postulates: states are complex vectors in a Hilbert space;
state evolution is governed by the Schrödinger equation; measurement
outcomes occur with probabilities given by the Born rule; physical
observables are represented by self-adjoint operators whose commutation
relations produce the uncertainty relations. Each postulate has been
verified to extraordinary empirical precision, and each plays a
well-defined role in the predictive machinery of the theory. At the
foundational level, however, no derivation is known within conventional
formulations that reduces any of the four postulates to consequences of
the others, and several foundational programs have attempted to supply
such a reduction from various directions. The Many-Worlds framework
attempts to derive the Born rule from decision-theoretic rationality
axioms applied to branching dynamics {[}wallace2012{]}. Gleason's
theorem derives the Born rule from assumptions of non-contextuality and
additivity on a Hilbert-space projector lattice, but the Hilbert-space
structure itself is assumed rather than derived {[}gleason1957{]}.
Bohmian mechanics reproduces the predictions of non-relativistic quantum
mechanics using particle positions and a non-local pilot wave
{[}bohm1952{]}. Relational and QBist programs reinterpret the postulates
at the interpretive level without modifying the formal content
{[}rovelli1996, fuchs2002{]}. None of these programs reduces all four
postulates to consequences of a single underlying structural framework.

This paper presents such a reduction within the Event Density (ED)
program. The program starts from a discrete ontology of events, chains,
participation relations, and commitment dynamics. A distinguished object
--- the participation measure --- is identified as the carrier of
quantum-mechanical content within this ontology. Under a clearly
specified thin-participation limit, the five-step derivation presented
in sections 3 through 7 produces each of the four quantum-mechanical
postulates as a structural consequence of the participation-measure
framework.

The argument is organized as follows. Section 2 describes the ED
primitive stack used in the derivation. Section 3 defines the
participation measure and establishes its structural properties. Section
4 derives the Schrödinger equation as the thin-participation limit of a
linear first-order complex evolution equation for the participation
measure. Section 5 derives the Born rule by showing that environmental
phase-randomization at commitment events, combined with an individuation
threshold, forces channel-selection probabilities to take the
bandwidth-fraction form, whose squared-amplitude structure is
definitional. Section 6 derives Bell-inequality-violating correlations
from non-factorizable joint participation measures and obtains the
Tsirelson bound from the sesquilinear inner product of the participation
manifold. Section 7 derives the Heisenberg uncertainty relations from
the Fourier-conjugate partition of the adjacency-band bandwidth budget.
Section 8 presents a structural observation --- the exponent-two thread
--- that connects the Born rule, the Schrödinger kinetic term, the
Madelung decomposition, the sublinear bandwidth-composition rule, and
the Fourier-uncertainty bound to a single definitional commitment.
Section 9 distinguishes structural content, which is derived from
primitives, from dimensional anchoring, which is inherited through the
ED Dimensional Atlas. Section 10 catalogues limitations and open
research frontiers, and section 11 concludes.

\subsection{2. The Event Density
Framework}\label{the-event-density-framework}

This section summarizes the ED primitive stack used in the derivation. A
complete account of the primitive stack is available in the companion
Event Density literature; the presentation here is restricted to the
structural content directly used in the derivations that follow.

The ED ontology is a graph-theoretic structure whose vertices are
micro-events and whose edges carry participation relations. A chain is a
sequence of micro-events held together by a consistent update rule; the
rule persists across the sequence and constitutes the chain's identity.
Participation is the relational substrate connecting micro-events; two
events participate when each integrates the other's becoming into its
own. Participation is not a force, field, or exchange of information at
the primitive level, but the relational fabric on which forces, fields,
and information-bearing structures are coarse-grained.

The scalar magnitude of a participation relation is called the
participation bandwidth. For a chain \(C\) at position \(x\) and time
\(t\), the total bandwidth budget \(b_{\mathrm{total}}(C)\) decomposes
into four distinguished bands: \begin{equation}
  b_K = b_K^{\mathrm{int}} + b_K^{\mathrm{adj}} + b_K^{\mathrm{env}} + b_K^{\mathrm{com}},
  \label{eq:fourband}
\end{equation} where \(b_K^{\mathrm{int}}\) denotes internal
rule-bandwidth sustaining the chain's own update rule,
\(b_K^{\mathrm{adj}}\) denotes adjacency bandwidth with the chain's
immediate participation-adjacent neighborhood, \(b_K^{\mathrm{env}}\)
denotes environmental bandwidth shared with the broader bath, and
\(b_K^{\mathrm{com}}\) denotes commitment-reserve bandwidth available
for commitment events. The index \(K\) ranges over the channels
available to the chain at the position and time of interest.

A channel is a stable bandwidth-preserving sub-structure of the
participation graph that supports a chain's update rule repeatedly
without rule-incompatibility. Channels are discrete at the graph level
and become a continuous index in the coarse-grained thick-regime
description. The multiplicity \(M\) of a chain at a position is the
count of distinct channels available for continuation. A commitment
event is the discrete selection of a single channel by a chain, together
with the addition of a new micro-event along the selected channel's
next-vertex position. The commitment event is local at the participation
site; its trigger condition is the rapid growth of environmental
bandwidth past an individuation threshold. Tension polarity is the phase
relation between the chain's update rule and the local flow direction;
it takes values in \(S^1\) and provides the phase content that the
participation bandwidth, a real non-negative scalar, does not carry on
its own.

The thin-participation regime is the limit in which the effective
multiplicity satisfies \(M_{\mathrm{eff}} \gg 1\), the environmental
band satisfies \(b_K^{\mathrm{env}} \to 0\), and the commitment rate
satisfies \(\Gamma_{\mathrm{commit}} \to 0\). In this regime the chain
is distributed coherently across many channels with no environmental
decoherence and no commitment events, constituting the primitive-level
analogue of an isolated quantum system.

\subsection{3. The Participation
Measure}\label{the-participation-measure}

The first step of the derivation is the definition of a single
mathematical object that carries both the amplitude content of the
participation bandwidth and the phase content of the tension polarity.

\subsubsection{3.1 Definition}\label{definition}

The participation measure is a complex-valued distribution over channels
and positions, defined by \begin{equation}
  P_K(x,t) = \sqrt{b_K(x,t)} \, e^{i \pi(K,x,t)},
  \label{eq:P_def}
\end{equation} where \(b_K \geq 0\) is the bandwidth in channel \(K\) at
position \(x\) and time \(t\), and \(\pi(K,x,t) \in [0, 2\pi)\) is the
tension-polarity phase. By construction, \begin{equation}
  b_K(x,t) = |P_K(x,t)|^2,
  \label{eq:b_is_Psquared}
\end{equation} so the participation bandwidth is the squared amplitude
of the participation measure. The event density, a scalar field that
counts micro-events per region, is recovered as the sum over channels:
\begin{equation}
  \rho(x,t) = \sum_K |P_K(x,t)|^2 = \sum_K b_K(x,t).
\end{equation} The coherent sum over channels defines a complex scalar
field \begin{equation}
  \Psi(x,t) = \sum_K P_K(x,t),
\end{equation} which plays the role of the quantum-mechanical
wavefunction in the thin-participation limit.

\subsubsection{3.2 Structural constraints}\label{structural-constraints}

The participation measure inherits several structural constraints from
the primitive stack. The measure must carry both amplitude and phase,
because the Born rule requires bandwidth weighting and interference
requires phase differences. The measure must support linear
superposition of the form \(\alpha P^A + \beta P^B\) for complex scalars
\(\alpha\), \(\beta\), because coherent recombination of channels at
interferometric junctions preserves the superposed structure. The
measure must admit a sesquilinear inner product of the form
\begin{equation}
  \langle P \, | \, Q \rangle = \sum_K \int dx \, P_K^*(x) \, Q_K(x),
  \label{eq:inner}
\end{equation} which yields a real non-negative norm equal to the total
participation bandwidth. The measure must be decomposable across the
four bands of equation\textasciitilde{}\eqref{eq:fourband}. Joint
participation measures on multipartite systems must admit
non-factorizable configurations, corresponding to entangled states.
Finally, the measure must be compatible with the commitment dynamics, in
which the bandwidth-fraction \(|P_{K^*}|^2 / \sum_K |P_K|^2\) provides
the selection probability at a commitment event.

\subsubsection{3.3 Uniqueness of the algebraic
structure}\label{uniqueness-of-the-algebraic-structure}

The participation measure must carry amplitude and phase, support linear
superposition with complex coefficients, and admit a sesquilinear inner
product. The minimum algebra satisfying these requirements is the
complex numbers \(\mathbb{C}\). Real scalars alone cannot support phase
rotations; the natural pair \((b, \pi)\) under custom operations is
isomorphic to \(\mathbb{C}\) via
equation\textasciitilde{}\eqref{eq:P_def}. Quaternions support all the
required operations but introduce non-commutative multiplication, which
conflicts with the symmetric structure of shared participation between
endpoints of an entangled pair. Higher-dimensional algebras are
over-structured relative to the primitive-level requirements. The
complex numbers constitute the unique minimum algebra for the
participation measure.

The square-root structure of equation\textasciitilde{}\eqref{eq:P_def}
is forced by the Madelung-class identification of the event density with
the squared modulus of the wavefunction in the quantum regime. Under
this identification, the bandwidth \(b_K\) must equal \(|P_K|^2\) rather
than \(|P_K|\), and the polar decomposition of a complex number fixes
the amplitude of \(P_K\) as \(\sqrt{b_K}\).

\subsection{4. The Schrödinger
Equation}\label{the-schruxf6dinger-equation}

The participation measure evolves in time according to a linear
first-order complex equation whose specific form is derivable from the
primitive stack together with the structural properties established in
section 3.

\subsubsection{4.1 Derivation of the evolution
equation}\label{derivation-of-the-evolution-equation}

The evolution of the participation measure is linear. Linearity follows
from the requirement that superposition be preserved under evolution,
which in turn follows from the multi-channel thin-regime structure and
the compatibility of the commitment dynamics with the Born rule. A
non-linear evolution would violate the preservation of superposition,
produce Born-rule inconsistency at successive commitment events, or
permit superluminal signaling in entangled systems; each of these
consequences is structurally excluded by the primitive-level
commitments.

The evolution is first-order in time. The chain's update rule applies
one step at a time; the rule's next application depends only on the
current state of the chain, not on past history. This is the
primitive-level content of a Markov-type dynamics, and it rules out
higher-order time derivatives.

The most general linear first-order evolution compatible with these
constraints is \begin{equation}
  i\hbar \, \partial_t P_K = H_K \, P_K + \sum_{K'} V_{KK'} \, P_{K'},
  \label{eq:P_evolution}
\end{equation} where \(H_K\) is a channel-diagonal term and \(V_{KK'}\)
represents off-diagonal inter-channel coupling. The factor \(i\hbar\) on
the left-hand side is forced by the following argument. Total bandwidth
of an isolated chain is conserved in the thin-participation regime, so
\begin{equation}
  \frac{d}{dt} \sum_K |P_K|^2 = 0.
\end{equation} Substituting
equation\textasciitilde{}\eqref{eq:P_evolution} and requiring this
vanishing for arbitrary \(P\) yields the anti-Hermitian condition
\(L_{KK'} + L_{K'K}^{*} = 0\) on the evolution kernel \(L\). Any
anti-Hermitian operator can be written as \(L = -iH/\hbar\) with \(H\)
Hermitian; the prefactor \(i\hbar\) then appears on the time-derivative
side, yielding equation\textasciitilde{}\eqref{eq:P_evolution}.

The numerical value of \(\hbar\) is inherited from the ED Dimensional
Atlas through the Madelung-class identification of the quantum regime
with the free-particle Schrödinger equation, whose diffusion coefficient
\(D = \hbar/(2m)\) is a mathematical identity under the Madelung
transformation. The structural form of
equation\textasciitilde{}\eqref{eq:P_evolution} is derived from
primitives; the specific numerical value of \(\hbar\) is inherited.

\subsubsection{4.2 Thin-participation
limit}\label{thin-participation-limit}

In the thin-participation limit the channel index \(K\) becomes
continuous. Under translation invariance of the participation graph, the
natural continuous index is the momentum \(k\); Stone's theorem on
unitary one-parameter groups {[}stone1930{]} identifies the generator of
translations as the momentum operator \(\hat{P} = -i\hbar \partial_x\).
In the momentum representation, the diagonal term of
equation\textasciitilde{}\eqref{eq:P_evolution} takes the form
\(H_k = \hbar^2 k^2 / (2m) + V\), where the quadratic dependence on
\(k\) is forced by rotation invariance, analyticity of the dispersion
relation at \(k=0\), and the non-relativistic regime; linear-in-\(k\)
and higher-order-in-\(k\) terms are excluded. The coefficient \(1/(2m)\)
follows from classical-correspondence (Ehrenfest) constraints, with the
mass parameter \(m\) inherited empirically at the dimensional-anchoring
level.

Summing equation\textasciitilde{}\eqref{eq:P_evolution} over channels
and identifying the result with \(\Psi = \sum_K P_K\) yields
\begin{equation}
  i\hbar \, \partial_t \Psi(x,t) = \left[-\frac{\hbar^2}{2m} \nabla^2 + V(x)\right] \Psi(x,t),
  \label{eq:Schrodinger}
\end{equation} the standard linear Schrödinger equation for a particle
of mass \(m\) in the potential \(V(x)\).

\subsubsection{4.3 Equivalence to the Madelung
form}\label{equivalence-to-the-madelung-form}

Polar decomposition of the wavefunction as
\(\Psi = \sqrt{\rho} \, e^{iS/\hbar}\), substituted into
equation\textasciitilde{}\eqref{eq:Schrodinger} and split into real and
imaginary parts, yields the continuity equation \begin{equation}
  \partial_t \rho + \nabla \cdot (\rho \mathbf{v}_s) = 0,
  \qquad \mathbf{v}_s = \frac{\nabla S}{m},
\end{equation} and the Hamilton-Jacobi equation with quantum-pressure
correction \begin{equation}
  \partial_t S + \frac{|\nabla S|^2}{2m} + V + Q = 0,
  \qquad Q = -\frac{\hbar^2}{2m} \frac{\nabla^2 \sqrt{\rho}}{\sqrt{\rho}}.
\end{equation} This is the Madelung form {[}madelung1927{]} of the
Schrödinger equation. The continuity equation and the
Hamilton-Jacobi-with-\(Q\) pair are mathematically equivalent to
equation\textasciitilde{}\eqref{eq:Schrodinger}.

\subsection{5. The Born Rule}\label{the-born-rule}

A measurement is a commitment event that occurs when environmental
bandwidth grows rapidly enough to force an individuation boundary
between the chain and its environment. At the commitment event, the
chain transitions from a multi-channel coherent state to a
single-channel state, and a new micro-event is added at the committed
channel's next-vertex position.

\subsubsection{5.1 Environmental
phase-randomization}\label{environmental-phase-randomization}

The environmental coupling that triggers a commitment event induces
random phase shifts on the chain's channels: \begin{equation}
  P_K(x,t) \;\longmapsto\; P_K(x,t) \, e^{i \delta_K(x,t)},
  \label{eq:phase_kick}
\end{equation} where \(\delta_K\) are random phases drawn from the
environmental-mode distribution. The phases \(\{\delta_K\}\) are
statistically independent across channels \(K\). Independence is forced
by three independent primitive-level arguments: the four-band
decomposition of equation\textasciitilde{}\eqref{eq:fourband} assigns
the environmental bandwidth of each channel to a distinct subset of
environmental modes; the locality of commitment events excludes
long-range environmental correlations between channels; and the
orthogonality of environmental modes, which are eigenmodes of a
Hermitian bath Hamiltonian, renders the phases of distinct modes
independent. Moreover, any correlation between \(\delta_K\) and
\(\delta_{K'}\) for \(K \neq K'\) would violate bandwidth conservation,
the locality of commitment, and the no-signaling property of entangled
bipartite systems; all three violations are structurally excluded.

\subsubsection{5.2 Destruction of inter-channel
coherence}\label{destruction-of-inter-channel-coherence}

Consider the squared magnitude of the coherent sum, \begin{equation}
  |\Psi|^2 = \sum_K |P_K|^2 + \sum_{K \neq K'} P_K^* \, P_{K'}.
\end{equation} Under the phase-kick of
equation\textasciitilde{}\eqref{eq:phase_kick}, the cross-terms acquire
random phase factors \(e^{i(\delta_{K'} - \delta_K)}\). Averaging over
the environmental phase distribution, \begin{equation}
  \langle e^{i(\delta_{K'} - \delta_K)} \rangle_{\mathrm{env}} = \delta_{K,K'}
\end{equation} for independent phases, so the cross-terms vanish on
environmental averaging. The surviving diagonal terms \(\sum_K |P_K|^2\)
constitute an incoherent mixture in the channel basis.

\subsubsection{5.3 Selection probability}\label{selection-probability}

The incoherent mixture assigns to channel \(K\) a weight \(|P_K|^2\).
Normalizing this weight by the total bandwidth \(\sum_K |P_K|^2\) yields
the selection probability \begin{equation}
  \mathrm{Prob}(K^*) = \frac{|P_{K^*}(x)|^2}{\sum_K |P_K(x)|^2}.
  \label{eq:Born}
\end{equation} Equation\textasciitilde{}\eqref{eq:Born} is the Born rule
{[}born1926{]}. Under the identification \(\Psi = \sum_K P_K\), the
amplitude in channel \(K\) equals the projection
\(\langle K | \Psi \rangle\), and
equation\textasciitilde{}\eqref{eq:Born} reads
\(\mathrm{Prob}(K^*) = |\langle K^* | \Psi \rangle|^2 / \langle \Psi | \Psi \rangle\),
recovering the conventional quantum-mechanical statement.

The squared-amplitude form of equation\textasciitilde{}\eqref{eq:Born}
is definitional rather than postulated. The identification
\(b_K = |P_K|^2\) from equation\textasciitilde{}\eqref{eq:b_is_Psquared}
makes the bandwidth-fraction and the squared-amplitude-fraction the same
quantity. The Born rule's exponent of two is not an independent
empirical fact requiring its own justification; it is a consequence of
the structural choice that constructs the bandwidth from the
participation measure by squaring the amplitude. Had the decomposition
been \(P_K = b_K e^{i\pi_K}\) with \(b_K = |P_K|\) instead, the Born
rule would read \(\mathrm{Prob}(K^*) = |P_{K^*}| / \sum_K |P_K|\). The
empirical consistency of the squared-amplitude form under Madelung
anchoring forces equation\textasciitilde{}\eqref{eq:b_is_Psquared}.

\subsection{6. Bell Correlations and the Tsirelson
Bound}\label{bell-correlations-and-the-tsirelson-bound}

Entanglement arises in the ED framework from shared participation
structure across two or more endpoints of a bipartite or multipartite
system. The joint participation measure of a bipartite system is a
complex-valued function \begin{equation}
  P^{AB}_{K_A, K_B}(x_A, x_B, t) \in \mathbb{C}
\end{equation} over the product channel space. A joint state is
factorizable if \(P^{AB}_{K_A K_B} = P^A_{K_A} \cdot P^B_{K_B}\), and
non-factorizable (entangled) otherwise.

\subsubsection{6.1 The CHSH scenario}\label{the-chsh-scenario}

Consider a bipartite system with subsystems \(A\) and \(B\) spatially
separated. Each subsystem admits two measurement settings, \(a\) and
\(a'\) for \(A\) and \(b\) and \(b'\) for \(B\), and each measurement
produces a binary outcome \(\pm 1\). The correlation functions
\begin{equation}
  E(a,b) = \sum_{s_A, s_B = \pm 1} s_A \, s_B \, \mathrm{Prob}(s_A, s_B \,|\, a, b)
\end{equation} combine into the CHSH statistic \begin{equation}
  S = E(a,b) + E(a,b') + E(a',b) - E(a',b').
  \label{eq:CHSH}
\end{equation} The Bell inequality {[}bell1964{]} states that
\(|S| \leq 2\) for any local hidden-variable theory; the CHSH
formulation {[}chsh1969{]} is the specific form of
equation\textasciitilde{}\eqref{eq:CHSH}. The Tsirelson bound
{[}tsirelson1980{]} states that \(|S| \leq 2\sqrt{2}\) under standard
quantum mechanics.

\subsubsection{6.2 Factorizable joint
participation}\label{factorizable-joint-participation}

For factorizable joint participation, the joint probability factors:
\(\mathrm{Prob}(K_A, K_B \,|\, a, b) = \mathrm{Prob}_A(K_A \,|\, a) \cdot \mathrm{Prob}_B(K_B \,|\, b)\).
Each subsystem's outcome depends on local measurements and local
randomness only. Standard Bell-CHSH algebra gives \begin{equation}
  S = \langle s_A(a) [s_B(b) + s_B(b')] \rangle + \langle s_A(a') [s_B(b) - s_B(b')] \rangle.
\end{equation} For realizations of the local random variables with
\(s_B(\cdot) \in \{-1, +1\}\), one of \(s_B(b) + s_B(b')\) and
\(s_B(b) - s_B(b')\) vanishes while the other has magnitude two.
Consequently \(|S| \leq 2\), which is the Bell-CHSH inequality.

\subsubsection{6.3 Non-factorizable joint
participation}\label{non-factorizable-joint-participation}

For non-factorizable joint participation, the joint probability does not
factor. Consider the prototypical singlet configuration \begin{equation}
  P^{AB}_{K_A K_B} = \frac{1}{\sqrt{2}} \bigl( \delta_{K_A,0}\delta_{K_B,1} - \delta_{K_A,1}\delta_{K_B,0} \bigr),
\end{equation} representing a pair of two-level subsystems in the
spin-singlet state. With measurement settings corresponding to Pauli
observables \(\sigma_a, \sigma_b\) on unit vectors
\(\mathbf{a}, \mathbf{b}\), the correlation function is \begin{equation}
  E(\mathbf{a}, \mathbf{b}) = -\mathbf{a} \cdot \mathbf{b}.
\end{equation} Choosing angles \(\mathbf{a} = \hat{\mathbf{x}}\),
\(\mathbf{a}' = \hat{\mathbf{z}}\),
\(\mathbf{b} = (\hat{\mathbf{x}} + \hat{\mathbf{z}})/\sqrt{2}\),
\(\mathbf{b}' = (\hat{\mathbf{x}} - \hat{\mathbf{z}})/\sqrt{2}\) yields
\(S = -2\sqrt{2}\), so \(|S| = 2\sqrt{2}\), saturating the Tsirelson
bound.

\subsubsection{6.4 The Tsirelson bound from participation-manifold
structure}\label{the-tsirelson-bound-from-participation-manifold-structure}

Define the CHSH operator \begin{equation}
  \hat{S} = \sigma_a \otimes \sigma_b + \sigma_a \otimes \sigma_{b'} + \sigma_{a'} \otimes \sigma_b - \sigma_{a'} \otimes \sigma_{b'}.
\end{equation} Its square satisfies \begin{equation}
  \hat{S}^2 = 4 \, I_A \otimes I_B + [\sigma_a, \sigma_{a'}] \otimes [\sigma_b, \sigma_{b'}],
\end{equation} using \(\sigma_a^2 = I\) for Hermitian unit-norm
observables. The operator-norm bound \(\|[\sigma, \sigma']\| \leq 2\)
for such observables gives \(\|\hat{S}^2\| \leq 4 + 2 \cdot 2 = 8\), so
\(\|\hat{S}\| \leq 2\sqrt{2}\). Since
\(|S| = |\langle P^{AB} | \hat{S} | P^{AB} \rangle| \leq \|\hat{S}\| \cdot \|P^{AB}\|^2\)
and the joint participation measure is normalized,
\(|S| \leq 2\sqrt{2}\). This Landau-Khalfin argument {[}landau1987{]}
rests on the sesquilinear inner-product structure of
equation\textasciitilde{}\eqref{eq:inner}, the normalization of the
joint participation measure, and the unit-norm Hermitian character of
the measurement observables, all of which are primitive-level or
U-derived properties of the participation-manifold.

The no-signaling property of the joint participation measure follows
from the locality of commitment events. The marginal probability on
\(A\) is independent of the measurement setting at \(B\); this is
structural rather than postulated.

\subsection{7. The Heisenberg Uncertainty
Relations}\label{the-heisenberg-uncertainty-relations}

The uncertainty relations arise in the ED framework from the
Fourier-conjugate decomposition of the adjacency-band bandwidth budget.
The adjacency bandwidth \(b^{\mathrm{adj}}\) encodes two structurally
distinct kinds of nearness: spatial nearness, associated with
localization in the position coordinate \(x\), and phase-propagation
nearness, associated with localization in the wavenumber \(k\) or
momentum \(p = \hbar k\). The two kinds of adjacency are connected by
the Fourier transform of the coherent-sum wavefunction.

\subsubsection{7.1 Adjacency-band
partition}\label{adjacency-band-partition}

The adjacency band admits a two-component orthogonal decomposition
\begin{equation}
  b^{\mathrm{adj}} = b_x + b_p,
  \label{eq:adj_partition}
\end{equation} where \(b_x\) is the position-adjacency bandwidth,
encoded in the spatial density distribution \(|\Psi(x)|^2\), and \(b_p\)
is the momentum-adjacency bandwidth, encoded in the momentum density
distribution \(|\tilde\Psi(p)|^2\). The two distributions are related by
the Fourier transform \begin{equation}
  \tilde\Psi(p,t) = (2\pi\hbar)^{-1/2} \int dx \, \Psi(x,t) \, e^{-ipx/\hbar}.
  \label{eq:Fourier}
\end{equation} Equation\textasciitilde{}\eqref{eq:Fourier} is the unique
basis-change between position and momentum representations under
translation invariance of the participation graph; Stone's theorem
identifies \(\hat{P} = -i\hbar \partial_x\) as the generator of
translations, and the Fourier transform is the basis-change to the
eigenbasis of \(\hat{P}\). Parseval's theorem
\(\int dx |\Psi|^2 = \int dp |\tilde\Psi|^2\) shows that the two
distributions are complementary decompositions of the same total
adjacency bandwidth. The Stone-von Neumann theorem establishes that the
Heisenberg-Weyl algebra generated by \(\hat{x}\) and \(\hat{P}\) is
irreducible on \(L^2(\mathbb{R})\), so no third independent
decomposition component exists in one spatial dimension.

\subsubsection{7.2 Derivation of the uncertainty
relation}\label{derivation-of-the-uncertainty-relation}

The bandwidth-allocation inequality for the orthogonal
partition\textasciitilde{}\eqref{eq:adj_partition} takes the form
\begin{equation}
  (\Delta b_x)(\Delta b_p) \geq K_{xp},
\end{equation} where \(K_{xp}\) is a structural constant. Under the
identifications \(b_x(x) \propto |\Psi(x)|^2\) and
\(b_p(p) \propto |\tilde\Psi(p)|^2\) inherited from the Born-rule
structure applied within each conjugate partition, the bandwidth spreads
are the standard position and momentum variances of the wavefunction,
\begin{equation}
  (\Delta b_x)^2 = \langle x^2 \rangle_{|\Psi|^2} - \langle x \rangle_{|\Psi|^2}^2 = (\Delta x)^2,
\end{equation} and similarly \((\Delta b_p)^2 = (\Delta p)^2\). The
mathematical Fourier-uncertainty theorem applied to the Fourier pair
\((\Psi, \tilde\Psi)\) with the canonical commutator
\([\hat{x}, \hat{P}] = i\hbar\) yields {[}heisenberg1927,
robertson1929{]} \begin{equation}
  \Delta x \cdot \Delta p \geq \frac{\hbar}{2}.
  \label{eq:Heisenberg}
\end{equation} The bound saturates for Gaussian minimum-uncertainty
wavepackets. Equation\textasciitilde{}\eqref{eq:Heisenberg} is the
Heisenberg position-momentum uncertainty relation. Unitary evolution
under equation\textasciitilde{}\eqref{eq:Schrodinger} preserves the
inequality at all times.

\subsubsection{7.3 Generalizations}\label{generalizations}

The four-band structure supports analogous decompositions for other
conjugate pairs. The internal rule-bandwidth and the commitment-reserve
bandwidth form a conjugate pair for which the energy-time uncertainty
relation holds. The commitment-reserve bandwidth and the rule-phase form
a conjugate pair for which the number-phase uncertainty relation holds.
Directional decompositions of the internal rule-bandwidth support the
angular-momentum uncertainty relations. In each case, the same
allocation-inequality mechanism within a specific band-pair of the
four-band structure produces the corresponding uncertainty relation. The
universal appearance of \(\hbar/2\) as the bound reflects the common
Fourier-conjugacy content, inherited through the anchoring of \(\hbar\)
at the dimensional level.

\subsection{8. The Exponent-Two Thread}\label{the-exponent-two-thread}

The five-step derivation of sections 3 through 7 produces the
quantum-mechanical content through a single structural commitment that
threads through otherwise-independent appearances of the number two.

\subsubsection{8.1 Five appearances}\label{five-appearances}

The Born rule selection probability
\(\mathrm{Prob}(K^*) = |P_{K^*}|^2 / \sum_K |P_K|^2\) contains the
squared amplitude. The Schrödinger kinetic term
\(\hat{H}_{\mathrm{kin}} = \hbar^2 k^2 / (2m)\) is quadratic in the
momentum. The Madelung decomposition
\(\Psi = \sqrt{\rho} \, e^{iS/\hbar}\) contains the square root of the
density. The sublinear bandwidth-composition rule expresses the combined
bandwidth of two channels in terms of squared-amplitude-level
cross-terms
\(b_{\mathrm{combined}}^2 = b_1^2 + b_2^2 + 2 c_{12} b_1 b_2\). The
Fourier-uncertainty bound \((\Delta x)^2 (\Delta p)^2 \geq (\hbar/2)^2\)
is quadratic in the variances. Each of these five appearances involves
the exponent two in a way that is not, in conventional quantum
mechanics, tied to a single underlying structural commitment.

\subsubsection{8.2 The single underlying
commitment}\label{the-single-underlying-commitment}

In the participation-measure framework, each of the five appearances
arises from the same definitional commitment: bandwidth is the squared
amplitude of the participation measure, \(b_K = |P_K|^2\)
(equation\textasciitilde{}\eqref{eq:b_is_Psquared}). The Born rule's
squared amplitude is the bandwidth fraction by construction. The
Schrödinger kinetic term's \(\hbar^2 k^2 / (2m)\) is the quadratic
leading-order dispersion forced by translation and rotation invariance;
the \(\hbar^2\) factor arises from the square of the momentum operator
acting on the wavefunction, which is the coherent sum of participation
amplitudes. The Madelung decomposition's \(\sqrt{\rho}\) is the
amplitude \(|\Psi|\) whose square is the density. The sublinear
composition rule's squared-amplitude cross-structure is the composition
law for the squared magnitudes of the participation measure. The
Fourier-uncertainty bound's squared variances arise from the
second-moment structure of the conjugate-partition densities. All five
are manifestations of the same structural choice.

In conventional quantum mechanics, these five appearances are logically
independent. The Born exponent could in principle differ from the
kinetic exponent without immediate empirical inconsistency; Gleason's
theorem and related arguments fix the Born exponent from Hilbert-space
axioms, but do not tie it to the kinetic exponent. In the
participation-measure framework, the five exponents are tied together
structurally: a change in the definitional identification \(b = |P|^2\)
would change all five. The structural unification is the program's
strongest positive claim.

\subsection{9. Structural Content versus Dimensional
Anchoring}\label{structural-content-versus-dimensional-anchoring}

The derivations of sections 3 through 7 produce the form of the
quantum-mechanical equations --- the Schrödinger structure, the Born
probability formula, the Bell-CHSH inequality and Tsirelson bound, and
the Heisenberg uncertainty relation --- from the primitive-level
participation-measure framework. The numerical values of the dimensional
constants that enter these equations are not derived from primitives.
The constant \(\hbar\) is anchored through the ED Dimensional Atlas via
the Madelung identification of the quantum regime's diffusion
coefficient \(D_{\mathrm{phys}} = \hbar/(2m)\) with the free-particle
Schrödinger equation; \(\hbar\) enters
equation\textasciitilde{}\eqref{eq:P_evolution} with its empirically
measured value. The mass parameter \(m\) is anchored empirically at the
species level; the primitive-level derivation of particle masses from
chain structure remains an open research problem and is flagged in
section 10.

The separation is clean at the level of the derivation. All structural
content --- linearity of the evolution equation, first-order-in-time
structure, operator-form decomposition of the evolution kernel, norm
preservation and the associated anti-Hermitian generator,
complex-valued-ness of the wavefunction, sesquilinear inner-product
structure, Fourier-conjugate partition of the adjacency band, the
squared-amplitude form of the Born rule, the Bell-CHSH factorizable
bound, the Tsirelson bound, the Heisenberg inequality --- is derived
from primitives combined with standard mathematical theorems (Stone's
theorem, Parseval's theorem, Stone-von Neumann theorem,
Fourier-uncertainty theorem, Landau-Khalfin operator-norm algebra). The
dimensional constants \(\hbar\) and \(m\) are inherited via the
Dimensional Atlas from empirical measurement. No other content is
inherited.

The framework does not predict the numerical value of \(\hbar\); it
inherits it. The framework does not predict the mass spectrum of
elementary particles; these are inherited at the species level and
constitute open research questions. The framework does not make
empirical predictions that differ from standard quantum mechanics within
the tested thin-participation regime; retrodiction attempts against
matter-wave interferometry data in the Eibenberger and Fein regimes
reproduce the standard predictions within the consistent-weak-form
epistemic category.

\subsection{10. Limitations and Research
Frontiers}\label{limitations-and-research-frontiers}

The structural reduction presented in this paper has explicit
limitations and leaves several research problems open for future work.

First, the primitive-level derivation of particle masses from chain
structure is flagged as speculative. The chain-mass problem asks how the
specific rule-type bandwidth signature of an electron, proton, or other
elementary particle is set by the underlying participation-graph
structure. A first-pass candidate identification of the rest mass with
the rule-type bandwidth divided by \(c^2\) has been articulated, but a
rigorous derivation requires a primitive-level derivation of the
rule-type taxonomy, which is itself an open question. The problem
connects to the origin of the Standard-Model mass spectrum and is a
multi-year research frontier.

Second, the numerical value of \(\hbar\) is inherited rather than
derived. The Dimensional Atlas anchors \(\hbar\) through the
Madelung-class identification of the quantum regime; a primitive-level
origin for the specific numerical value of \(\hbar\) would reduce the
inheritance further but has not been attempted.

Third, the derivation is non-relativistic. The relativistic extension
would replace the quadratic dispersion \(H_k = \hbar^2 k^2 / (2m)\) with
the relativistic dispersion \(E = \sqrt{p^2 c^2 + m^2 c^4}\), requiring
a Lorentz-covariant formulation of the participation measure and a
corresponding covariant four-band decomposition. The Klein-Gordon and
Dirac equations would arise as the thin-limit consequences of the
covariant framework; neither has been derived in the present work.

Fourth, the derivation is single-particle. A multi-particle or
quantum-field-theoretic extension would require the development of
second-quantized participation measures, with particle-creation and
particle-annihilation operators emerging as primitive-level structural
consequences. The connection to the Standard Model field content would
require the rule-type taxonomy derivation mentioned above.

Fifth, the evolution equation\textasciitilde{}\eqref{eq:P_evolution} has
been derived in the isolated-chain thin-participation limit, where
environmental coupling and commitment events are suppressed. The
extension to non-isolated chains would produce a
Lindblad-master-equation-type structure with jump operators
corresponding to environmental commitment events; the primitive-level
derivation of this extended dynamics is near-term future work.

Sixth, the framework does not currently make empirical predictions that
distinguish it from standard quantum mechanics in the regime where
standard quantum mechanics has been tested. Distinguishing content is
expected to emerge in regimes where the thin-participation limit breaks
down, including the quantum-to-classical transition at macroscopic scale
and cosmological-scale dynamics relevant to baryogenesis; preliminary
retrodiction attempts in these regimes have been structurally
informative but not empirically distinguishing within currently
available data.

\subsection{11. Conclusion}\label{conclusion}

The non-relativistic quantum mechanics of a single particle has been
shown to arise as a structural consequence of the Event Density
primitive stack. A single mathematical object --- the participation
measure \(P_K(x,t) = \sqrt{b_K(x,t)} \, e^{i\pi(K,x,t)}\) --- serves as
the carrier of quantum-mechanical content. Its thin-participation limit
produces the Schrödinger equation; commitment events under environmental
phase-randomization produce the Born rule in its squared-amplitude form;
non-factorizable joint participation measures produce
Bell-inequality-violating correlations saturating at the Tsirelson
bound; Fourier-conjugate partition of the adjacency bandwidth produces
the Heisenberg uncertainty relations. The five otherwise-independent
appearances of the exponent two in quantum mechanics --- the Born rule,
the Schrödinger kinetic term, the Madelung decomposition, the sublinear
bandwidth-composition rule, and the Fourier-uncertainty bound --- are
manifestations of a single definitional commitment, the identification
of bandwidth with squared amplitude. No additional physical postulate
beyond the participation-measure structure, the four-band bandwidth
decomposition, and the commitment dynamics is required. The numerical
values of \(\hbar\) and \(m\) are inherited through the ED Dimensional
Atlas; the primitive-level derivation of particle masses remains an open
research problem. The result is a structural reduction of the four
independent postulates of standard quantum mechanics --- state
structure, unitary evolution, measurement probability, and
observable-commutator uncertainty --- to structural consequences of a
single primitive-level framework.

\begin{center}\rule{0.5\linewidth}{0.5pt}\end{center}

\subsection{References}\label{references}

{[}bell1964{]} Bell, J. S. (1964). On the Einstein-Podolsky-Rosen
paradox. \emph{Physics Physique Fizika}, \textbf{1}, 195--200.

{[}bohm1952{]} Bohm, D. (1952). A suggested interpretation of the
quantum theory in terms of ``hidden'' variables. I and II.
\emph{Physical Review}, \textbf{85}, 166--179, 180--193.

{[}born1926{]} Born, M. (1926). Zur Quantenmechanik der Stoßvorgänge.
\emph{Zeitschrift für Physik}, \textbf{37}, 863--867.

{[}chsh1969{]} Clauser, J. F., Horne, M. A., Shimony, A., and Holt, R.
A. (1969). Proposed experiment to test local hidden-variable theories.
\emph{Physical Review Letters}, \textbf{23}, 880--884.

{[}fuchs2002{]} Fuchs, C. A. (2002). Quantum mechanics as quantum
information (and only a little more). \emph{arXiv:quant-ph/0205039}.

{[}gleason1957{]} Gleason, A. M. (1957). Measures on the closed
subspaces of a Hilbert space. \emph{Journal of Mathematics and
Mechanics}, \textbf{6}, 885--893.

{[}heisenberg1927{]} Heisenberg, W. (1927). Über den anschaulichen
Inhalt der quantentheoretischen Kinematik und Mechanik.
\emph{Zeitschrift für Physik}, \textbf{43}, 172--198.

{[}landau1987{]} Landau, L. J. (1987). On the violation of Bell's
inequality in quantum theory. \emph{Physics Letters A}, \textbf{120},
54--56.

{[}madelung1927{]} Madelung, E. (1927). Quantentheorie in
hydrodynamischer Form. \emph{Zeitschrift für Physik}, \textbf{40},
322--326.

{[}robertson1929{]} Robertson, H. P. (1929). The uncertainty principle.
\emph{Physical Review}, \textbf{34}, 163--164.

{[}rovelli1996{]} Rovelli, C. (1996). Relational quantum mechanics.
\emph{International Journal of Theoretical Physics}, \textbf{35},
1637--1678.

{[}stone1930{]} Stone, M. H. (1930). Linear transformations in Hilbert
space, III: operational methods and group theory. \emph{Proceedings of
the National Academy of Sciences}, \textbf{16}, 172--175.

{[}tsirelson1980{]} Tsirelson, B. S. (1980). Quantum generalizations of
Bell's inequality. \emph{Letters in Mathematical Physics}, \textbf{4},
93--100.

{[}wallace2012{]} Wallace, D. (2012). \emph{The Emergent Multiverse:
Quantum Theory according to the Everett Interpretation}. Oxford
University Press.

\end{document}
